\section{Method}
Optimisation experiments involved both the practical methods associated with
 preparation and characterisation of thermoset composite samples as well as
 computational techniques developed to carry out Bayesian optimisation based
 on data obtained.
Furthermore, synthetic routines were developed for producing any materials
 not readily available commercially.
Initially, this section will present practical methods including
 synthesis of itaconic acid polyester and diesters, thermoset composite
 preparation, and characterisation techniques employed.
Subsequently, computational method by which Bayesian optimisation was
 implemented will be elucidated.

% ========================================================================
\subsection{Materials}
The propylene glycol (99~\unit{\percent}), itaconic acid
 (99~\unit{\percent}), anhydrous calcium sulphate (99~\unit{\percent}), and
  hydroquinone (99~\unit{\percent}) were supplied by ``Thermo Scientific 
  Chemicals'', a division of \textit{Thermo Fisher Scientific Inc.} based 
  in Waltham, Massachusetts, U.S.
The toluene (99.8~\unit{\percent}), dried magnesium sulphate 
 (99~\unit{\percent}), sodium chloride (99.5~\unit{\percent}), and 
 dichloromethane (\>95~\unit{\percent}) were supplied by ``Fisher Chemicals'',
 a division of \textit{Thermo Fisher Scientific Inc.}
The mixed xylenes (97~\unit{\percent}) were supplied by ``Alfa Aesar'' a 
 division of \textit{Thermo Fisher Scientific Inc.}
The methanol (99.9~\unit{\percent}) was supplied by ``Honeywell Riedel-de 
 Ha\"{e}n'', a brand under \textit{Honeywell International Inc.} based in
 Charlotte, North Carolina, U.S.
The p-toluenesulphonic acid monohydrate (98~\unit{\percent}) and 
 chloroform-d (99.8~\unit{\percent}) were supplied by ``Sigma-Aldrich'', a
 subsidiary of \textit{Merck KGaA} based in Darmstadt, Germany.
Dibenzoyl peroxide (20~\unit{\percent}) filled with calcium sulphate was 
 supplied by \textit{Nouryon Specialty Chemicals B.V.} based in Amsterdam,
 The Netherlands.
Sodium bicarbonate (99~\unit{\percent}) was supplied by 
 \textit{Fluorochem Ltd} based in Glossop, Derbyshire, U.K.
 Phenolphthalein solution (1~\unit{\percent} w/w in ethanol) and potassium
 hydroxide (0.56~\unit{\percent} w/w in methanol) were supplied 
 by \textit{Scientific Laboratory Supplies Ltd} based in the East Riding 
 of Yorkshire, U.K.

% ========================================================================
\subsection{Background}
This research centred on the optimisation of a simple bio-based composite
 comprised of a thermosetting polymer matrix and a particulate reinforcing
 material.

\paragraph{Thermoset Polymer Matrix}
The thermosetting matrix was based on a cross-linked unsaturated polyester
 system; these rely upon a combination of condensation and addition polymer.
The condensation polymer is typically an unsaturated linear polyester
 formed from the copolymerisation of diols and diacids, with at least one
 monomer having an unsaturated double bond in its structure which is
 retained after copolymerisation- this copolymer is referred to as an
 unsaturated polyester \cite{Seraji2022}.
Whilst the addition polymer is typically generated through a radical-addition
 homopolymerisation reaction, where a monomer which had been acting as
 an unsaturated polyester solvent reacts and crosslinks individual 
 unsaturated polyester chains- these monomers are therefore referred to as
 reactive diluents \cite{Chen2015}.
The combination of unsaturated polyester and reactive diluent is known
 as an unsaturated polyester resin; a viscous liquid, which, on addition
 of radical-initiating chemicals, can rapidly form a rigid thermoset
 material.

Previous work by Panic et al. found dimethyl itaconate to be the strongest
 contender of the diitaconates as a bio-based replacement for styrene
 in unsaturated polyester resins, given that resultant cured thermosets
 exhibited relatively higher crosslink densities (73~\unit{\percent} the
 density of a styrene-based thermoset) \cite{Panic2017}.
This work therefore sought to replicate these unsaturated polyester resins,
 where the reactive diluent was pure dimethyl itaconate and the unsaturated
 polyester was poly(propylene itaconate).

\paragraph{Particulate Reinforcing}
Previous work by Kharas et al. made use of particulate calcium sulphate
 in the preparation of bioresorbable bone cement composites, where it's 
 content was known to effect the final thermosets' strength and modulus
 \cite{Kharas1997}.
This work therefore employed the readily-available anhydrous
 calcium sulphate in preparation as the particulate reinforcing element
 of the bio-based composites considered.

% ========================================================================
\subsection{Synthesis of Diitaconate}
Dimethyl itaconate was synthesised by esterification of itaconic acid
 with methanol by the Fischer method which had seen previous success
 in the literature \cite{Trejbal2022,Cowie1979,Cowie1977}; this is
 presented in Figure~\ref{fig:SimpleDMIReactionScheme} as a simplified
 reaction scheme.
Method was adapted from the work of Trejbal et al. and involved the mixing
 of itaconic acid in an excess of methanol in a molar ratio of 1:3
 \cite{Trejbal2022}.
P-toluenesulphonic acid was used as a heterogenous strong acid catalyst
 and made up 4~\unit{\percent}~w/w of the final reaction mixture.
\begin{figure*}[htb]
    \begin{center}
        \captionsetup{format=custom,width=0.8\textwidth,skip=0pt}
        \includegraphics[width=0.9\textwidth]{SimpleDMIReactionScheme}
        \captionof{figure}{
            \footnotesize{
                Simple reaction scheme for synthesis of dimethyl itaconate.
            }
        }
        \label{fig:SimpleDMIReactionScheme}
    \end{center}
\end{figure*}
Figure~\ref{fig:SynthesisRigDMI_Lettered} shows how this reaction mixture
 was placed in the lower round-bottomed flask
 (Fig.\ref{fig:SynthesisRigDMI_Lettered}D), and brought to a  vigorous boil
 at around 130~\unit{\degreeCelsius}.
Vapours of both water and methanol passed up an insulated glass tube
 (Fig.\ref{fig:SynthesisRigDMI_Lettered}E) and were caught in a distillation
 column (Fig.\ref{fig:SynthesisRigDMI_Lettered}G) which was connected to
 a reboiler (Fig.\ref{fig:SynthesisRigDMI_Lettered}F) at it's base.
Temperature at the column's top (Fig.\ref{fig:SynthesisRigDMI_Lettered}H)
 was retained at around 70~\unit{\degreeCelsius} by adjusting the temperature of the
 reboiler (120~\unit{\degreeCelsius}) and amount of insulation applied around the column;
 this ensured a steady stream of methanol-rich vapours passing through the
 Claisen head (Fig.\ref{fig:SynthesisRigDMI_Lettered}I) and re-condensing
 within the Liebig condenser (Fig.\ref{fig:SynthesisRigDMI_Lettered}J).
Methanol-rich liquid was therewith returned to the lower round-bottomed
 flask via a dropping funnel (Fig.\ref{fig:SynthesisRigDMI_Lettered}K),
 whilst water was caught in the reboiler
 (Fig.\ref{fig:SynthesisRigDMI_Lettered}F), pushing the reversible
 esterification reaction toward the desired diester product.

Before initiating the reaction by starting both the hotplate stirrers
 (Fig.\ref{fig:SynthesisRigDMI_Lettered}A), an argon-rich atmosphere was
 generated by using a Schlenk line to cycle a vacuum (500~\unit{\milli\bar})
 and argon gas (350~\unit{\milli\litre\per\minute}); this was to discourage
 unwanted oxidation reactions and reduce the chances of vapour ignition.
During the reaction however a valve was opened to allow pressure 
 equilibration between the reaction vessels and the atmosphere; mitigating
 the risk of explosion.
Methanol cycling was halted 3.5~\unit{\hour} into the reaction by closing the
 dropping funnel's tap (Fig.\ref{fig:SynthesisRigDMI_Lettered}K), after
 which, all heating was ceased on reaching 4~\unit{\hour}.
\noindent
\begin{minipage}{0.485\textwidth}
    \includegraphics[width=\textwidth]{SynthesisRigDMI_Lettered.png}
    \captionsetup{format=custom,width=\textwidth,skip=0pt}
    \captionof{figure}{
        \footnotesize{
            Apparatus for synthesis of dimethyl itaconate \cite{Chemix2025}.
            (A) Hotplate stirrer
            (B) Sand bath
            (C) Temperature controller probe
            (D) Three-necked round-bottomed flask containing the reaction
                mixture
            (E) Insulated glass tube
            (F) Single-necked round-bottomed flask acting as a distillation
                column reboiler containing water
            (G) Distillation column
            (H) Thermometer
            (I) Four-way Claisen head
            (J) Liebig condenser
            (K) Dropping funnel for recycled methanol
        }
    }
    \label{fig:SynthesisRigDMI_Lettered}
\end{minipage}
Workup was carried out in toluene; trial-and-error experimentation found
 this to be suitable for extraction of dimethyl itaconate.
Final product mixture retrieved from the round-bottomed flask
 (Fig.\ref{fig:SynthesisRigDMI_Lettered}D) was placed in a separating funnel
 with toluene in a 1:3 volume ratio.
Separating funnel contents was washed with water thrice, saturated
 sodium bicarbonate solution twice, and saturated sodium chloride once;
 each washing step involved a volume ratio of 2:1 between the previously
 obtained organic fraction and the next volume of washing liquid used.
A gravity filtration step was carried out between the water and sodium
 bicarbonate washing steps to remove insolubles.
Following washing, the organic fraction was removed to an Erlenmeyer flask
 and dried using anhydrous magnesium sulphate; thereafter removed using
 gravity filtration.
Dry organics were placed in a rotary evaporator and heated at
 40~\unit{\degreeCelsius} whilst the pressure was dropped from an initial
 70~\unit{\milli\bar} to 20~\unit{\milli\bar} over an average period of
 8~\unit{\hour}; this removed all but a small mass of toluene and water
 from the liquid product.
From this liquid, dimethyl itaconate was crystallised out at room
 temperature by addition of a nucleating agent.
The nucleating agent was crystalline dimethyl itaconate; serendipitously,
 storage of a particularly pure batch of rotary evaporation product in a
 fridge at 3~\unit{\degreeCelsius} lead to the spontaneous nucleation
 required for freezing- generating the first crystalline sample of dimethyl
 itaconate.
The above method afforded a 45~\unit{\percent} yield although higher values
 may be possible if an extended period of boiling were adopted.

% ========================================================================
\subsection{Diitaconate Characterisation}
Characterisation techniques were used to monitor the reaction progress,
 identify chemicals, and ascertain the product purity; these included
 thin layer chromatography, Fourier-transform infrared spectroscopy, and
 nuclear magnetic resonance spectroscopy.

Initial impressions of purity were garnered from spotting chemicals onto
 thin layer chromatography plates (polyester with a layer of silica gel and
 ultraviolet indicator; 6~\unit{\milli\metre} pore diameter and
 0.2~\unit{\milli\metre} layer thickness) and using a mixture of
 dichloromethane and methanol (2.5:7.5 v/v respectively) as the mobile phase.
On completion, spot(s) could be observed under ultraviolet light
 (254~\unit{\nano\metre} wavelength) and their retardation factors
 (\(R_{F}\)) calculated based on equation \ref{eq:RetardationFactor}
 relating this metric with the distance migrated by the substance in
 question (\(b\) in \unit{\milli\metre}) and that travelled by the
 mobile phase (\(a\) in \unit{\milli\metre}) \cite{Ettre1993}.
\begin{equation}\label{eq:RetardationFactor}
    R_{F}=\frac{b}{a}
\end{equation}
Having solvated and diluted pure itaconic acid to a 1.25~\unit{\percent}~w/w
 solution (using stock mobile phase), it was spotted on a plate- this
 produced a single spot with an \(R_{F}\) value of 0.79.
Meanwhile, having spotted a 5~\unit{\percent}~w/w solution of final
 crystalline product (expected to be dimethyl itaconate) a single spot with
 an \(R_{F}\) of 0.91 was obtained.
These results provided benchmarks by which reaction progress could be
 gauged as well as supporting the hypothesis that these two crystalline
 materials were generally quite pure, although we could not rule out the
 presence of non-fluorescent or trace chemicals using this technique.
Samples taken periodically from the reaction vessel
 (Fig.\ref{fig:SynthesisRigDMI_Lettered}D) during the course of an
 esterification reaction revealed a change in composition: the first sample
 (retrieved after 15\unit{\minute} of heating) presented a single spot with
 an \(R_{F}\) value of 0.81, whilst a final sample (retrieved after 
 4~\unit{\hour}) presented two spots with \(R_{F}\) values of 0.90 and 0.81
 respectively.
Given the similarity of these samples with those of the crystalline
 reactant and product analysed previously, these results indicate the
 partial conversion of the reactant itaconic acid into the desired product
 dimethyl itaconate.

Vibrational spectra were obtained using a Nicolet iS50 Fourier-transform
 infrared spectrometer using the attenuated total reflectance technique;
 the spectral range between 400-4000~\unit{\per\centi\metre} was
 investigated at a resolution of 0.09~\unit{\per\centi\metre}.
Figure \ref{fig:FTIRComparison} presents the spectra obtained for the pure
 reactant (itaconic acid) and compares it with that of the product
 isolated using the rotary evaporator; this diagram supports the hypothesis
 that the product was succesfully esterified to form the product dimethyl
 itaconate in a number of ways:
Firstly, the broad peak characteristic of O-H bond stretching commonly
 located between 2500-3200~\unit{\per\centi\metre} is clearly
 present for the itaconic acid but not for the product
 (Fig.\ref{fig:FTIRComparison}A), suggesting the loss of O-H bonds
 and therewith the carboxyl groups associated with either the initial
 reactant (itaconic acid) or the intermediate product (monoethyl itaconate).
Secondly, strong peaks are observed around 1710~\unit{\per\centi\metre}
 (Fig.\ref{fig:FTIRComparison}B\&\ref{fig:FTIRComparison}D) which
 is commonly indicative of C=O bond stretching (which occurs between 
 1640-1760~\unit{\per\centi\metre}); this would be consistent with C=O
 bond retention which would occur during succesful esterification.

 Furthermore, it may be noted that peaks associated with carboxylic C=O bonds
 occur between 1640-1730~\unit{\per\centi\metre} whilst ester C=O peaks
 may occur between 1690-1750\textsuperscript{-1}; the positive difference
 between the reactant C=O peak located at 1700~\unit{\per\centi\metre}
 (Fig.\ref{fig:FTIRComparison}B) and the product's peak at
 1720~\unit{\per\centi\metre} (Fig.\ref{fig:FTIRComparison}D) would
 support succesful esterification.
Thirdly, an intermediate-sized peak is observed both for the reactant
 (Fig.\ref{fig:FTIRComparison}C) and the product
 (Fig.\ref{fig:FTIRComparison}E) at 1630 and
 1640~\unit{\per\centi\metre} respectively; their size and location are
 suggestive of nonconjugated C=C bond stretching (occuring between
 1620-1680~\unit{\per\centi\metre}), which would be consistent with
 retention of the double bond during the proposed esterification.
As an aside, clusters of peaks likely associated with C-H bond stretching
 were observed on both spectra
 (Fig.\ref{fig:FTIRComparison}F\&\ref{fig:FTIRComparison}G) but
 inter-comparison was obfuscated by the strong and broad O-H stretching
 signal apparent in the reactant's spectrum.
\noindent
\begin{minipage}{0.485\textwidth}
    \includegraphics[width=\textwidth]{FTIRComparison.png}
    \captionsetup{format=custom,width=\textwidth,skip=0pt}
    \captionof{figure}{
        \footnotesize{
            FT-IR spectra of reactant and product.
            (A) 2500-3200~\unit{\per\centi\metre}.
            (B) 1700~\unit{\per\centi\metre}.
            (C) 1630~\unit{\per\centi\metre}.
            (D) 1720~\unit{\per\centi\metre}.
            (E) 1640~\unit{\per\centi\metre}.
            (F) 2950 \& 3050~\unit{\per\centi\metre}.
            (G) 2875, 2980 \& 3035~\unit{\per\centi\metre}.
        }
    }
    \label{fig:FTIRComparison}
\end{minipage}
Both \textsuperscript{1}H and \textsuperscript{13}C nuclear magnetic
 resonance spectroscopy was carried out on a sample of product using a
 ``Bruker Avance III HD 500~\unit{\mega\hertz} Smart Probe'' spectrometer,
 with spectrograms produced at 500~\unit{\mega\hertz} and and
 125~\unit{\mega\hertz} respectively.
Chemical shifts are considered as downfield shifts (\(\delta\) in ppm) from
 the tetramethyl silane standard (\(\delta=0\)~ppm) and calibrated using the
 shift exhibited by trace amounts of non-deuterated chloroform
 (\(\delta=7.26\)~ppm singlet for \textsuperscript{1}H and \(\delta=77\)~ppm 
 triplet for \textsuperscript{13}C) within the sample solvent.
The chloroform-d solvent also introduced a trace amount of water which
 exhibited a singlet peak (\(\delta =1.63\)~ppm) in the \textsuperscript{1}H
 spectrogram; these peaks were ignored when examining the spectra for
 dimethyl itaconate, as presented here:

Dimethyl itaconate: \textsuperscript{1}H NMR (CDCl\textsubscript{3},
 500~\unit{\mega\hertz}): \(\delta\) 3.34 (d, 2H, \(J=0.71\)~Hz), 3.69
 (s, 3H), 3.76 (s, 3H), 5.71 (q, 1H, \(J=0.97\)~Hz), 6.32 (s, 1H).
 \textsuperscript{13}C \{\textsuperscript{1}H\} NMR (CDCl\textsubscript{3},
 125~\unit{\mega\hertz}): \(\delta\) 37.5, 52.1, 52.2, 128.6, 133.7, 166.6,
 171.1.
[GO BACK OVER THIS WITH THE AID OF "AN OLD CAPTIVITY NOTES FROM 22ND
JANUARY 2025"]

A quantitative \textsuperscript{1}H spectrogram showed water constituting a
 5~\unit{\percent}~w/w impurity, although the dimethyl itaconate sample
 purity is expected to exceed 95~\unit{\percent} given that the deuterated
 solvent will have introduced an unknown amount of water.

% ========================================================================
\subsection{Synthesis of Unsaturated Polyester Resin}
Unsaturated polyester was synthesised via the polyesterification of
 itaconic acid with propylene glycol to produce poly(propylene itaconate),
 as described by the simplified reaction scheme in
 Figure~\ref{fig:SimpleUPReactionScheme}.
This was achieved using the melt polycondensation method since a precedent
 exists in the literature for its succesful use in the synthesis of
 linear unsaturated polyesters derived from itaconic acid
 \cite{Dai2017,Spasojevic2021,Pantic2022,Panic2017}.
The stoichiometric ratio between itaconic acid and propylene glycol was
 set at 1:1.1, on the basis that propylene glycol is lost during the
 heating process which reaches temperatures exceeding it's boiling point
 \cite{Parkyn1967}.
\begin{figure*}[htb]
    \begin{center}
        \captionsetup{format=custom,width=0.8\textwidth,skip=0pt}
        \includegraphics[width=0.9\textwidth]{SimpleUPReactionScheme.png}
        \captionof{figure}{
            \footnotesize{
                Simple reaction scheme for synthesis of poly(propylene itaconate).
            }
        }
        \label{fig:SimpleUPReactionScheme}
    \end{center}
\end{figure*}
Figure~\ref{fig:SynthesisRigUP_Lettered} shows the apparatus used to
 synthesise unsaturated polyester; note, a Dean-Stark trap
 (Fig.\ref{fig:SynthesisRigUP_Lettered}I) has been used allowing the
 `azeotropic method' of polyester synthesis (as opposed to
 the `dry' method) which reduces the time taken in synthesising a
 high molecular weight polyester \cite{Parkyn1967}.
In this method, an entraining agent (mixed xylenes) which can form an
 azeotrope with the water is added to the mixture (5~\unit{\percent} w/w)
 \cite{Korbar1993,Bodnar1990}; this increases
 the rate at which water is condensed and trapped by the Dean-Stark,
 driving the reaction forward by removing the unwanted byproduct of water
 from the reaction mixture.
Furthermore, argon gas was used to sparge the reaction mixture
 (Fig.\ref{fig:SynthesisRigUP_Lettered}E) facilitating the removal of
 water from the bulk mixture and mitigating the combustion risk associated
 with high temperature xylenes \cite{Parkyn1967}.

The procedure was initiated by charging the four-necked round-bottomed
 flask (Fig.\ref{fig:SynthesisRigUP_Lettered}D) with the requisite masses
 of the two monomers and xylenes, opening the argon taps to provide a
 5 bubbles per second of flow of sparging gas through the argon
 injector (Fig.\ref{fig:SynthesisRigUP_Lettered}E) \cite{Parkyn1967},
 opening the water taps to actively cool the Liebig condenser
 (Fig.\ref{fig:SynthesisRigUP_Lettered}I), setting the overhead stirrer
 to a speed of 50~rpm (Fig.\ref{fig:SynthesisRigUP_Lettered}G)
 \cite{Parkyn1967}, and bringing the reaction mixtures temperature
 (Fig.\ref{fig:SynthesisRigUP_Lettered}F) up to 160~\unit{degreeCelsius}
 using the hotplate (Fig.\ref{fig:SynthesisRigUP_Lettered}A).
During the procedure the mixture's temperature was raised monotonically 
 by 10~\unit{\degreeCelsius\per\hour} until it reached a maximum of
 210~\unit{\degreeCelsius}, at which it was held.
\noindent
\begin{minipage}{0.485\textwidth}
    \includegraphics[width=\textwidth]{SynthesisRigUP_Lettered.png}
    \captionsetup{format=custom,width=\textwidth,skip=0pt}
    \captionof{figure}{
        \footnotesize{
            Apparatus for synthesis of unsaturated poly(propyl itaconate) \cite{Chemix2025}.
            (A) Hotplate stirrer
            (B) Sand bath
            (C) Temperature controller probe
            (D) Four-necked round-bottomed flask containing the reaction
                mixture
            (E) Stainless-steel argon injector
            (F) Probes for batch and vapour tamperatures
            (G) Overhead stirrer
            (H) Dean-Stark trap
            (G) Liebig condenser.
        }
    }
    \label{fig:SynthesisRigUP_Lettered}
\end{minipage}
Reaction progress was ascertained by periodically retrieving a sample of the
 reaction mixture and ascertaining it's acid number, whose unit is
 \unit{\milli\gram\per\gram} describing the milligrams of potassium hydroxide
 required to neutralise a gram of sampled unsaturated polyester.
This rough form of end-group analysis is commonly used in unsaturated
 polyester synthesis \cite{Pantic2022,Spasojevic2021,Costa2017,Parkyn1967}
 and provides a proxy value estimating the number of un-esterified acid
 groups remaining within the unsaturated polyester.
An American Society for Testing and Materials standard (``ASTM D4662-20'')
 was used for this purpose.
In-line with common practice, heating of the polyesterifying mixture was
 ceased at an acid value of 50 \cite{Parkyn1967,Bodnar1990,Korbar1993}.
On cessation of heating, an outlet was opened for egress of remnant xylenes
 and the unsaturated polyester was allowed to cool in an air bath.

Since cooling unsaturated polyester will undergo a monotonic increase in
 viscosity until it reaches room temperature where it exists as a
 solid thermoplastic mass; it is therefore common to prepare the unsaturated
 polyester resin immediately after cessation of heating by adding
 reactive diluent directly to the cooling unsaturated polyester
 \cite{Parkyn1967}.
Unfortunately, reactive diluents, by their very nature, may be liable to
 polymerise during this addition causing a premature gelation (untimely
 and incomplete cross-linking of unsaturated polyester resin usually
 culminating in the genesis of a rubber-like solid).
Mitigation against the dual-threat of hardening into thermoplastic and
 gelation into rubber is commonly enacted by awaiting a suitable intermediate
 temperature and addition of an addition-polymerisation inhibiting agent
 \cite{Parkyn1967,Ahamad2001,Cherian2004}.
In this work, we awaited a temperature of 100~\unit{\degreeCelsius} before 
 adding the requisite masses of inhibitor and reactive diluent.
This work adopted the commonly-used free-radical scavenger hydroquinone
 as an inhibiting agent, and this was added so as to make up
 0.06~\unit{\percent}~w/w of the cooling unsaturated polyester mixture;
 hydroquinone was dissolved in a small amount of propylene glycol before
 addition to ensure homogeneity during it's mixing into the bulk
 \cite{Parkyn1967}.
The inhibitor percentage was determined through a process of
 trial-and-error, where lower percentages tended to facilitate premature
 gelation of unsaturated polyester resins on the shelf, whilst higher
 percentages caused problems during the curing procedures.
Note, the percentage of inhibiting agent was above that which is
 commonly used in the synthesis of conventional unsaturated polyester
 resins (approximately 0.02~\unit{\percent}) \cite{Parkyn1967,Cherian2004};
 this was the consequence of both the relatively higher proportion of
 unsaturated sites along the polyester backbone and a higher copolymerisation
 affinity between the reactive diluent and the sites of unsaturation.
Subsequent of introduction and homogenisation of inhibiting agent into
 the unsaturated polyester mixture, the reactive diluent (dimethyl itaconate)
 was added such that it made up 30~\unit{\percent}~w/w of the final
 unsaturated polyester resin.
Reactive diluent content was intentionally less than the commonly used
 optimal ratios seen in literature, as this was a characteristic that was
 to be optimised for \cite{Parkyn1967,Ahamad2001,Cherian2004,Bodnar1990}.

% ========================================================================
\subsection{Unsaturated Polyester Resin Characterisation}
[WRITE ABOUT GPC OR NMR IF COMPLETED]

% ========================================================================
\subsection{Composite Preparation}
Figure \ref{fig:CompositePreparation} presents an overview of the method by
 which cured thermoset composite samples of various compositions were
 prepared.
Given the danger of self-accelerating decomposition reactions and associated
 thermal runaway, an external manufacturer was responsible for the mixing
 and provision of a safely-phlegmatised base initiating agent comprised of
 dibenzoyl peroxide and calcium sulphate
 (Fig.\ref{fig:CompositePreparation}F); this arrived pre-mixed in a set ratio
 of 2:8 between the dibenzoyl peroxide (\(1-c_{2}=0.2\) w/w) and the
 calcium sulphate (\(c_{2}=0.8\) w/w).
Base unsaturated polyester resin was mixed
 (Fig.\ref{fig:CompositePreparation}G) in a set ratio of 3:7 between
 unsaturated polyester resin (\(c_{1}=0.7\) w/w) and dimethyl
 itaconate (\(1-c_{1}=0.3\) w/w) as described previously; the elected
 ratio represents the lowest amount of reactive diluent possible before
 viscosity precluded the possibility of further mixing- this increased
 the size of the possible parameter space to be interrogated (by increasing
 the range of modified unsaturated polyester resin reactive diluent
 concentrations that could be mixed to by means of dilution).

Modified unsaturated polyester resin was mixed
 (Fig.\ref{fig:CompositePreparation}E) in a varying ratio of
 \(x_{1}:(1-x_{1})\) between the base unsaturated polyester resin
 (\(x_{1}\) w/w) and the dimethyl itaconate (\(1-x_{1}\) w/w); by associating
 the \(x_{1}\) parameter with a range (\([0.45,1.00]\)) within which it
 could vary, we fully describe the first predictor variable to be optimised
 (\(x_{1}\), unsaturated polyester content modifier).
Modified initiating agent was mixed (Fig.\ref{fig:CompositePreparation}D)
 in a varying ratio of \(x_{2}:(1-x_{2})\) between the base initiating agent
 (\(x_{2}\)) and calcium sulphate (\(1-x_{2}\)); \(x_{2}\) was provided
 a range (\([0.4,1.00]\)) hence describing the second predictor variable
 to be optimised (\(x_{2}\), dibenzoyl peroxide content modifier).
Initiated unsaturated polyester resin was mixed
 (Fig.\ref{fig:CompositePreparation}C) in a varying ratio of
 \(x_{3}:(1-x_{3})\) between modified unsaturated polyester resin
 (\(x_{3}\)) and modified initiating agent (\(1-x_{3}\)); \(x_{3}\) was 
 provided a range (\([0.5,0.95]\)) which therewith described the final
 predictor to be optimised (\(x_{3}\), initiating agent content modifier).

Initiated unsaturated polyester resin samples were thermochemically cured
 into rigid thermoset plastic samples by exposing them to elevated
 temperatures in a drying oven (Fig.\ref{fig:CompositePreparation}B); the
 curing profile was fixed for all samples: 3~hours at 
 80~\unit{\degreeCelsius} followed by 1~\unit{\hour} at
 120~\unit{\degreeCelsius}.
All cured thermoset samples were machined to the same dimensions
 (Fig.\ref{fig:CompositePreparation}A) in preparation for characterisation;
 this was achieved by a grinding process.

\noindent
\begin{minipage}{0.485\textwidth}
    \includegraphics[width=\textwidth]{CompositePreparation.png}
    \captionsetup{format=custom,width=\textwidth,skip=0pt}
    \captionof{figure}{
        \footnotesize{
            Mixing and processing schedule for the manufacture of thermoset
            composite samples.
            (A) Machining.
            (B) Thermochemical curing.
            (C) Mixing for initiated sample.
            (D) Mixing for modified initiator.
            (E) Mixing for modified resin.
            (F) Manufacturer's mixing for base initiator.
            (G) In-house mixing for base resin.
        }
    }
    \label{fig:CompositePreparation}
\end{minipage}
% ========================================================================
\subsection{Composite Characterisation}
Given that a key aspect of a thermoset's utility lies in it's strength and
 stiffness within the elastic region of deformation we considered the yield
 strength and Young's modulus as legitimate objective variables of interest.
Compressive testing was elected as a means by which these metrics could be
 obtained from samples based on both the simplicity with which a curing
 regime could be devised for the manufacture of samples of the requisite
 dimensions and the sparing mass of material required for these.
An American Society for Testing and Materials standard was employed 
 (``ASTM D695-23'') in the devising of method; with the cylindrical test
 sample form, sample diameter of 12.7~\unit{\milli\metre}, and the
 compression speed of 1.3~\unit{\milli\metre\per\minute} adopted.
However, sample length was adapted down by half (to 12.7~\unit{\milli\metre})
 to reduce the mass of material required per sample.

\paragraph{Equipment}
Compressive testing was carried out using a ``Tinius Olsen 50ST''
 universal test machine (50~\unit{\kilo\newton} capacity) fitted with
 compression platens.
Standard software (``Horizon'') purveyed by the machine's manufacturer
 was used in the control of the machine, retrieval of data from the load
 cells, and export of formatted data for further analysis.
Exported data consisted of time (\(\pm\)~0.16~\unit{\micro\second}),
 force (\(\pm\)~1~\unit{\milli\newton}), and platen displacement
 (\(\pm\)~1~\unit{\nano\metre}).

\paragraph{Euler Buckling}
To mitigate against concerns of Euler buckling taking place in samples of
 the proposed dimensions, we calculated the critical load based on
 Euler's critical load equation and some estimated properties for
 conventional polyester thermosets \cite{Euler1744}.
Equation \ref{eq:EulerCriticalLoad} shows how the critical load 
 (\(P_{cr}\) in \unit{\newton}) is related to the elastic modulus of the
 material (\(E\) in \unit{\pascal}), the second moment of area (\(I\) in
 \unit{\metre^{4}}),
 the effective length factor (\(K\), unitless), and the initial length of
 the column in question (\(L\) in \unit{\metre}).
\begin{equation}\label{eq:EulerCriticalLoad}
    P_{cr}=\frac{\pi EI}{(KL)^2}
\end{equation}
Young's moduli for conventional thermoset polyesters were expected to range
 from 2-4.5~\unit{\giga\pascal} \cite{Lawrence1960}.
The second moment of the area for the top surface of the cylindrical sample
 with radius 6.35~\unit{\milli\metre} was calculated to be
 \(1.276\times10^{-9}\)~\unit{\metre^{4}} by assuming it to be a perfect
 filled circular area, for which the generic equation
 \ref{eq:CircularSecondMomentOfArea} applied (where the \(r\) equated to the
 radius of the circle in \unit{\metre}).
\begin{equation}\label{eq:CircularSecondMomentOfArea}
    I_{y}=\frac{\pi}{4}r^{4}
\end{equation}
By assuming that both ends of the column would be braced against rotation
 or translation (given their compression between two platens), the
 theoretical column effective length factor would have been 0.5; as a
 further precaution a conservative design value of 0.65 was assumed
 based on pertinent literature \cite{AISC1989}.
On these assumptions it was estimated that the samples in question would
 be likely to Euler buckle between 433-976~\unit{\kilo\newton}.
As the compressive strength expected of conventional unsaturated
 polyester thermosets were to be expected to fall within the range of
 89-252~\unit{\mega\pascal} and this corresponded to maximum compressive
 forces of between 11-32~\unit{\kilo\newton} (given the surface area of
 the samples), we were reassured that the samples were unlikely to
 experience Euler buckling.

\paragraph{Correction of Raw Data}
Imperfect rigidity of the test machine itself confounded the raw exported
 datasets generated for samples tested.
By carrying out a compressive test with no sample between the platens,
 a force-displacement curve representing the inherent stiffness of the
 machine was generated- this was the correctional dataset.
Equation \ref{eq:CorrectedDisplacement} shows how the corrected displacement
 value (\(d\) in \unit{\milli\metre}) could be obtained by taking a
 raw displacement value (\(d_{raw}\) in \unit{\milli\metre}) and
 subtraacting the corresponding correctional displacement value
 (\(rd_{cor}\)).
\begin{equation}\label{eq:CorrectedDisplacement}
    d = d_{raw} - d_{cor}
\end{equation}\

Figure \ref{fig:CorrectionalMethod} presents the correctional dataset, an
 archetypical raw dataset, and it's corresponding corrected dataset.
Point~`A' is corrected to produce point~`B' by using it's associated force
 to find \(d_{cor}\), after which this can be subtracted from it's
 \(d_{raw}\) value to produce the \(d\) value that point~`b' obtains.
\noindent
\begin{minipage}{0.485\textwidth}
    \includegraphics[width=\textwidth]{CorrectionalMethod.jpg}
    \captionsetup{format=custom,width=\textwidth,skip=0pt}
    \captionof{figure}{
        \footnotesize{
            Strain correction of an archetypical raw dataset; the legend
            notes the various datasets considered.
        }
    }
    \label{fig:CorrectionalMethod}
\end{minipage}

\paragraph{Processing of Corrected Data}
Force values from a sample's corrected dataset were processed into
 engineering stress values by taking into account the initial surface area.
Each cylindrical sample's diameter was measured using a ``Magnusson''
 digital caliper (accuracy \(\pm\)~0.02~\unit{\milli\metre}) from three
 different positions (\(\text{ø}_{1},\text{ø}_{2},\text{ø}_{3}\)
 in \unit{\milli\metre}), each offset by roughly +45\degree~around the
 circumference, and the mean of these values (\(\text{ø}_{avg}\) in
 \unit{\milli\metre}) calculated (see Equation
 \ref{eq:AverageTopSurfaceDiameter}).
\begin{equation}\label{eq:AverageTopSurfaceDiameter}
    \text{ø}_{avg} = \frac{\text{ø}_{1}+\text{ø}_{2}+\text{ø}_{3}}{3}
\end{equation}
Equation \ref{eq:AverageTopSurfaceArea} shows how the initial top-face
 surface area was calculated from the average top surface diameter
 \(\text{ø}_{avg}\).
 \begin{equation}\label{eq:AverageTopSurfaceArea}
    A_{init} = 0.25 * \pi * \text{ø}_{avg}^{2}
\end{equation}
Subsequently, engineering stress (\(\sigma\) in \unit{\mega\pascal}) could
 be calculated (Equation \ref{eq:EngineeringStress}) by considering the
 force applied (\(F\) in \unit{\newton}) and the initial top surface area,
 \(A_{init}\) \cite{Martin2006}.
\begin{equation}\label{eq:EngineeringStress}
    \sigma _{eng} =\frac{F}{A_{init}}
\end{equation}

Displacement values from a sample's corrected dataset were processed into
 engineering strain values by taking into account the initial length of the
 samples and the length of the sample during compression.
Initial length (\(l_{init}\) in \unit{\milli\metre}) was calculated
 (\ref{eq:AverageInitialSampleLength}) by taking the mean of three length
 measurements (\(l_{1},l_{2},l_{3}\) in \unit{\milli\metre}) (these taken in
 the same fashion as for diameter).
 \begin{equation}\label{eq:AverageInitialSampleLength}
    l_{init} = \frac{l_{1}+l_{2}+l_{3}}{3}
\end{equation}
Equation~\ref{eq:SampleLengthUnderCompression} shows the simple link between
 length of a sample at a given point during a compressive test (\(l\) in
 \unit{\milli\metre}), the initial sample length (\(l_{init}\)) and the
 corrected displacement at that point \(d\).
\begin{equation}\label{eq:SampleLengthUnderCompression}
    l = l_{init}-d
\end{equation}
Subsequently, engineering strain was calculated (\ref{eq:EngineeringStrain})
 by finding the absolute change in sample length (\(\Delta l_{init}\) in 
 \unit{\milli\metre}) and dividing this by the initial length of the
 sample (\(l_{init}\)).
\begin{equation}\label{eq:EngineeringStrain}
    \epsilon _{eng}=\frac{\Delta l_{init}}{l_{init}}=\frac{l-l_{init}}{l_{init}}
\end{equation}

\paragraph{Stress-Strain Smoothing}
Sample-rate reduction and filtering were applied simultaneously during the
 downsampling of the stress-strain dataset resulting in a smoothed dataset.
Downsampling was achieved using a windowed averaging technique, whereby
 all the points within the unsmoothed dataset as a whole (\(p_{us}\)) was
 split (Eq.\ref{eq:WindowDivisionOfDataset}) into a number of windows along
 the strain-axis (\(p_{s}\)) by electing a suitable number of points per
 window (\(p_{w}\)).
Suitability was determined by trial-and-error within the wider experimental
 framework but was also based on the requirement of a perfect divisibility
 of \(p_{us}\) by \(p_{w}\).
Given the constant size of the unsmoothed dataset (\(p_{us}=1,000\)) a
 constant number of points per window was elected (\(p_{w}=5\)).
\begin{equation}\label{eq:WindowDivisionOfDataset}
    p_{s}=p_{us}/p_{w}
\end{equation}
Equation~\ref{eq:WindowAveraging} shows how the processs of averaging stress
 (\(\sigma_{eng}^{i}\)) and strain (\(\epsilon_{eng}^{i}\)) values within
 a given window yielded a single point described by it's average stress
 (\(\overline{\sigma}_{eng}^{j}\)) and strain (\(\overline{\epsilon}_{eng}^{j}\)).
\begin{equation}\label{eq:WindowAveraging}
    \overline{\sigma}_{eng}^{j},\overline{\epsilon}_{eng}^{j}=\frac{\sum_{i=1}^{p_{w}}\sigma_{eng}^{i}}{p_{w}},\frac{\sum_{i=1}^{p_{w}}\epsilon_{eng}^{i}}{p_{w}}
\end{equation}
Averaged points from each window were assembled to form the new downsampled
 dataset depicted in Equation~\ref{eq:SmoothedDataset}.
\footnotesize
\begin{equation}\label{eq:SmoothedDataset}
    [[\overline{\sigma}_{eng}^{j},\overline{\epsilon}_{eng}^{j}],[\overline{\sigma}_{eng}^{j+1},\overline{\epsilon}_{eng}^{j+1}],...,[\overline{\sigma}_{eng}^{p_{s}},\overline{\epsilon}_{eng}^{p_{s}}]]
\end{equation}
\normalsize
Figure~\ref{fig:AverageSmoothing} demonstrates the efficacy of this
 technique in removing noise from the unsmoothed stress-strain dataset.
\noindent
\begin{minipage}{0.485\textwidth}
    \includegraphics[width=\textwidth]{AverageSmoothing.jpg}
    \captionsetup{format=custom,width=\textwidth,skip=0pt}
    \captionof{figure}{
        \footnotesize{
            Multi-scale depiction of the stress-strain dataset before
            and after downscaling.
            (A) Overview line plot.
            (B) Zoomed scatter plot.
        }
    }
    \label{fig:AverageSmoothing}
\end{minipage}

\paragraph{Derivative Analysis}
The proportional limit may be defined as the stress at which a material
 fails to abide by Hooke's law (e.g. stress fails to proportionally rise
 with strain)- the gradient of the preceding elastic region of deformation
 being equal to the Young's modulus.
Given the material's unreliability in producing stress-strain curves with
 an obvious linear region, the limit of proportionality was equated to the
 point at which the stress-strain curve began to level-off (i.e. the first
 point of inflection).
To find this inflection point, the first derivative of stress with respect
 to strain (tangent modulus \(E_{t}\) in \unit{\mega\pascal}) was calculated
 from the smoothed stress-strain dataset (Eq.\ref{eq:TangentModulus}) by
 taking into account the instantaneous change in both stress (\(\Delta
 \sigma_{eng}\)) and strain (\(\Delta \epsilon_{eng}\)).

\begin{equation}\label{eq:TangentModulus}
    E_{t}=\frac{\Delta \sigma_{eng}}{\Delta \epsilon_{eng}}=\frac{\sigma_{eng}^{j+1}-\sigma_{eng}^{j}}{\epsilon_{eng}^{j+1}-\epsilon_{eng}^{j}}
\end{equation}

Figure~\ref{fig:DerivativeComparison} shows how smoothing of the 
 stress-strain dataset facilitated the analysis of the resultant tangent
 modulus versus strain dataset.
Note that markedly step-like behaviour in the raw stress-strain data
 (Fig.\ref{fig:DerivativeComparison}) was an artifact of an unknown
 processing step applied by the universal test machine's proprietary
 software.

\noindent
\begin{minipage}{0.485\textwidth}
    \includegraphics[width=\textwidth]{DerivativeComparison.jpg}
    \captionsetup{format=custom,width=\textwidth,skip=0pt}
    \captionof{figure}{
        \footnotesize{
            A comparison of the tangent modulus vs strain curves obtained
            from both the raw and the smoothed counterpart stress-strain
            datasets.
            (A) Strain at first inflection point.
        }
    }
    \label{fig:DerivativeComparison}
\end{minipage}

Detection of the first inflection point was automated using a simple
 algorithm that iterated chronologically through the tangent modulus versus
 strain dataset from low strains to high.
At each point the algorithm would calculate an average gradient exhibited
 by an equal number of points both ahead and behind (the 'point horizon')
 the current strain value; cessation and return of the current strain was
 enacted on reaching a point at which summed gradient ahead and behind
 fell below a threshold of zero.
Figure~\ref{fig:DerivativeComparison}A shows how a point horizon of 10 was
 used to identify the first inflection point, and therewith the strain at
 which the limit of proportionality had been reached.

\paragraph{Objective Retrieval}
By re-considering the smoothed stress-strain dataset, the point with the
 closest proximity to the proportional limit's determined strain value was
 equated to the first inflection point: the proportional limit
 (Fig.\ref{fig:ExampleStressStrain}A).
\noindent
\begin{minipage}{0.485\textwidth}
    \includegraphics[width=\textwidth]{ExampleStressStrain.jpg}
    \captionsetup{format=custom,width=\textwidth,skip=0pt}
    \captionof{figure}{
        \footnotesize{
            Means by which Young's modulus and yield point were acquired
            from an archetypical stress-strain dataset.
            (A) First inflection point.
            (B) Second inflection point.
        }
    }
    \label{fig:ExampleStressStrain}
\end{minipage}

As shown in Figure~\ref{fig:ExampleStressStrain}, a polynomial was fit to
 the ten points preceding the proportional limit; allowing approximation of
 Young's modulus (\(E\) in \unit{\mega\pascal}) from the change in stress
 (Fig.\ref{fig:ExampleStressStrain}: \(\Delta \sigma _{eng}^{A}\)) and strain
 (Fig\ref{fig:ExampleStressStrain}: \(\Delta \epsilon _{eng}^{A}\)) in
 accordance with Equation~\ref{eq:YoungsModulusCalculation}
 \cite{Horwood2023}.
Outputted Young's moduli were negative given the compresive nature
 of the dataset, but were naturally taken as absolute values.
 
\begin{equation}\label{eq:YoungsModulusCalculation}
    E=\frac{\sigma_{eng}}{\epsilon_{eng}}=\frac{\Delta \sigma _{eng}^{A}}{\Delta \epsilon _{eng}^{A}}
\end{equation}

Based on it's common definition (``the stress at which a material begins to
 behave in a plastic fashion''), this work defined
 (Eq.\ref{eq:YieldStrengthDefinition}) the yield strength (\(\sigma_{y}\)
 in \unit{\mega\pascal}) as the stress required to reach the second
 inflection point (Fig.\ref{fig:ExampleStressStrain}B) as visualised
 by \(\sigma_{eng}^{B}\) in Figure~\ref{fig:ExampleStressStrain}
 \cite{Vaidya2018}.

\begin{equation}\label{eq:YieldStrengthDefinition}
    \sigma_{y}=\sigma_{eng}^{B}
\end{equation}

Since the samples tested consistently produced similar stress-strain curves
 to that depicted in Figure~\ref{fig:ExampleStressStrain}, the second
 inflection point was trivially located searching the stress-strain dataset
 for the point exhibiting the highest stress value; the sample's yield or
 break strength.

 % ========================================================================
\subsection{Optimisation}
As described above, three predictors controlled the composition of samples:
 the unsaturated polyester content ratio (\(x_{1}\)), the dibenzoyl peroxide
 content modifier (\(x_{2}\)) and the initiating agent content modifier
 (\(x_{3}\)).
These three predictors and their ranges formed a d-dimensional (\(d=3\))
 parameter space (\(\mathcal{X}\)) which is a subset of the entire real coordinate
 space (\(\mathcal{X}\subseteq \mathbb{R}^{d}\)).
This parameter space was explored by the sequential optimisation techniques,
 with the objective variables including Young's modulus (\(E\)) and yield
 strength (\(\sigma_{y}\)).

Three different optimisation procedures were carried out from the same
 initial prior dataset: two were single-objective Bayesian routines with
 differing objectives, whilst the third was a purely exploratory method
 which was objective-indifferent.
A fourth multi-objective Bayesian optimisation routine incorporating both
 the objectives previously considered was carried out using a combination of
 the original prior dataset and all data generated by the three sequential
 optimisations.

\paragraph{Single-Objective Experiments}
Optimisation problems were based on black-box functions \(f\) which
 took any point within the parameter space \(\mathcal{X}\) and returned a real-valued
 target \(\mathbb{R}\); this being a mapping between the two respective
 spaces (\(f:\mathcal{X}\rightarrow \mathbb{R}\)).
Equation~\ref{eq:GoverningOptimisation} describes optimisation as the
 maximisation of the black-box function by tuning the predictor variable
 values (\(\mathbf{x}=[x_{1},x_{2},x_{3}]\)) such that the global optimum
 set of predictors (\(\mathbf{x}^{*}=[x_{1}^{*},x_{2}^{*},x_{3}^{*}]\)) is
 sought.

\begin{equation}\label{eq:GoverningOptimisation}
    \mathbf{x}^{*}=\argmaxA_{\mathbf{x}\in\mathcal{X}} f([x_{1},x_{2},x_{3}])
\end{equation}

Equation~\ref{eq:SOBOProblems} describes the two black-box functions which
 were separately maximised during the course of the single-objective
 optimisations; the former \(f_{E}\) and latter \(f_{\sigma}\) considering
 Young's modulus \(E\) and yield strength \(\sigma_{y}\) as target variables
 respectively.

\begin{equation}\label{eq:SOBOProblems}
    \begin{split}
        &E=f_{E}(\mathbf{x})\\
        &\sigma_{y}=f_{\sigma}(\mathbf{x})
    \end{split}
\end{equation}

Eight vectors of predictor variables, each corresponding to a sample (\(i\)),
 were representatively drawn by Sobol sampling from across the parameter
 space \(\mathcal{X}\) forming an initial matrix of predictors (\(\mathbf{X}\)); the
 matrix's form is presented in Equation~\ref{eq:MatrixOfPredictors}.

\begin{equation}\label{eq:MatrixOfPredictors}
    \begin{split}
        \mathbf{X}=&[\mathbf{x}^{i=1},\mathbf{x}^{i=2},...,\mathbf{x}^{i=8}]\\
        =&[[x_{1}^{i=1},x_{2}^{i=1},x_{3}^{i=1}],[x_{1}^{i=2},...],...]
    \end{split}
\end{equation}

Evaluation of all the vectors of the predictor matrix (\(f_{E}(\mathbf{X})\)
 or \(f_{\sigma}(\mathbf{X})\)) generated a vector of targets (\(\mathbf{E}\) or
 \(\bm{\sigma_{y}}\)).
Combination of the initial predictor matrix \(\mathbf{X}\) and the
 corresponding targets \(f(\mathbf{X})\) produced an initial dataset
 (\(\mathcal{D}=[\mathbf{X},f(\mathbf{X})]\)) with which sequential
 optimisation experiments were begun.
During the course of a sequential optimisation experiment, 16 further samples
 were drawn from across the parameter space \(\mathcal{X}\), with one sample taken per
 iteration; each sample's predictor vector \(\mathbf{x}\) and function
 evaluation \(f(\mathbf{x})\) was appended to the growing dataset
 \(\mathcal{D}\) associated with that experiment.
At each iteration the Bayesian experiments strove to balance exploration and
 exploitation during selection of sample composition within the parameter
 space, whilst the benchmark deterministic experiment solely focussed upon
 exploration.

Bayesian optimisation involved an iterative process of developing a surrogate
 model approximating the latent objective function based on the available
 dataset \(\mathcal{D}\), using an acquisition function to evaluate
 candidate suitability and selection of a new sample \(\mathbf{x}^{i}\), and
 finally evaluation of a sample followed by incorporation of the resultant
 target scalar \(\sigma_{y}^{i}\) and predictor vector \(\mathbf{x}^{i}\)
 into the growing dataset \(\mathcal{D}\).

On updating the dataset \(\mathcal{D}\), Gaussian process regression was used
 to develop a Gaussian process regression model \(\mathcal{GP}\) which can
 be considered a distribution over possible functions
 (\(f(\mathbf{x})\sim \mathcal{GP}(m(\mathbf{x}),K(\mathbf{x},\mathbf{x}'))\))
 fully described by both it's mean (\(m\)) and covariance (\(K\)) functions.
Computationally, the Gaussian process was modelled by a high-dimensional
 multivariate Gaussian distribution \(\mathcal{N}(\bm{\mu},\bm{\Sigma})\)
 described by it's mean vector \(\bm{\mu}\) and covariance matrix
 \(\bm{\Sigma}\); each succesive iteration saw the prior process conditioned
 on new data producing an updated posterior process.

For the purpose of encoding our prior beliefs in process smoothness we
 employed the Mat\'{e}rn family of covariance functions; specifically,
 one employing \(\nu=\)\nicefrac{5}{2} which yields
 Equation~\ref{eq:MaternCovarianceFunction} based on Euclidean distance
 (\(d\)).
The \(\nu=\nicefrac{5}{2}\) case balanced the need for smoothness in sampled
 paths with the drawbacks of higher values of \(\nu\); as
 \(\nu \rightarrow \infty\), the sampled paths become infinitely
 differentiable, making them unrealistically smooth for practical
 application \cite{Snoek2012} and potentially interfering with succesful
 calculation of the likelihood \cite{Stein1999}.
\footnotesize
\begin{equation}\label{eq:MaternCovarianceFunction}
    \begin{split}
        K_{M\nicefrac{5}{2}}&(x,x')=(1+\sqrt{5}d+\frac{5}{3}d^{2})exp(-\sqrt{5}d)\\
        &\text{where } d=|x-x'|
    \end{split}
\end{equation}
\normalsize

Each iteration saw the use of a one-step lookahead sampling policy utilising
 a simple reward utility function (Equation~\ref{eq:UtilityFunction}) which
 takes the current dataset \(\mathcal{D}\) and of all samples \(\mathbf{X}\)
 returns the sample \(\mathbf{x}\) exhibiting the highest target
 (\(\phi\); in practice either \(E\) or \(\sigma_{y}\)) according to the
 posterior process's mean function \(max\ m(X)\).
\begin{equation}\label{eq:UtilityFunction}
    u(\mathcal{D})=max\ m(\mathbf{X})
\end{equation}

Utility quantification of a putative sample \(\mathbf{x}_{p}\) was
 calculated by the ``expected improvement'' acquisition function
 \(\alpha_{EI}\) (Equation~\ref{eq:ExpectedImprovement}) which calculates
 the marginal gain in utility (\(u(\mathbb{D'})-u(\mathcal{D})\)) between
 the incumbent best target value in the current dataset (\(\phi^{*}\)) and
 that of a new dataset containing the putative \(\mathbf{x}_{p}\) value
 (\(u(\mathbb{D'}\))).
The second line of Equation~\ref{eq:ExpectedImprovement} clarifies how the
 calculation is one for the expected marginal gain, whilst the third
 presents reformulation for the noiseless case; this equation exemplifies
 how putative points are rewarded based on both scale and probability of
 improvement effectively trading off exploitation and exploration.
This work practically applied a more sophisticated logged noise-aware
 implementation (``qLogNoisyExpectedImprovement'' or ``qLNEI'') which has
 recently seen succesful use in high-dimensional computational domains
 \cite{Ament2025}.
\footnotesize
\begin{equation}\label{eq:ExpectedImprovement}
    \begin{split}
        &\alpha_{EI}(\mathbf{x}_{p},\mathcal{D})=u(\mathbb{D'})-u(\mathcal{D})\\
        &=\mathbb{E}[max\,m_{\mathcal{D}'}(\mathbf{X}')]-max\,m_{\mathcal{D}}(\mathbf{X})\\
        &=\int_{\phi^{*}}^{\infty}\phi\mathcal{N}(\mu,\sigma^{2})d\phi-\phi^{*}\int_{\phi^{*}}^{\infty}\mathcal{N}(\mu,\sigma^{2})d\phi\\
    \end{split}
\end{equation}
\normalsize

Trivial optimisation of the acquisition function
(Equation~\ref{eq:AcquisitionFunctionOptimisation}) was achieved by
 generating a large set of Sobol sampled points across the parameter space
 \(\mathbf{X}_{p}\), after which cheap evaluation facilitated identification of
 the candidate sample with the maximal utility score \(\mathbf{x}^{+}\).
\begin{equation}\label{eq:AcquisitionFunctionOptimisation}
    \mathbf{x}^{+}=\argmaxA_{\mathbf{X}_{p}\in\mathcal{X}}\alpha(\mathbf{X}_{p})
\end{equation}

The objective-indifferent space-filling benchmark sampling exercise was based
 on a sequential Monte Carlo strategy by the name of ``mc-intersite-proj-th''
 (``MIPT'') \cite{Crombecq2011}.
Utility was quantified by the conditional acquisition function displayed
 in Equation~\ref{eq:MIPT}, which took both a putative point
 \(\mathbf{x}_{p}\) and the current dataset's previous sample coordinates
 \(\mathbf{X}\) and returns either zero or the Euclidean distance to the
 nearest previously sampled point \(k\) depending on whether projected
 distance is below or in excess of the \(d_{m}\) value respectively.
The \(d_{m}\) is dependent upon the number of previous samples \(n\) and the
 arbitrary constant \(s\) which has been found to be optimally set to 0.5
 for a reasonable trade-off between intersite distance and projected
 distance.
\footnotesize
\begin{equation}\label{eq:MIPT}
    \begin{split}
        &\alpha_{mipt}(\mathbf{x}_{p},\mathbf{X}) =\\
        &\begin{cases}
            0       & if\:\min\limits_{\forall \mathbf{x}_i \in X}||\mathbf{x_{i}-\mathbf{x_{p}}}||_{-\infty}<d_{m}\\
            k,      & if\:\min\limits_{\forall \mathbf{x}_i \in X}||\mathbf{x_{i}-\mathbf{x_{p}}}||_{-\infty}\ge d_{m}
         \end{cases}\\
        &\text{where}\:d_{m}=\frac{2s}{n}\:\text{and}\:k=\min\limits_{\forall \mathbf{x}_i \in X}||\mathbf{x_{i}-\mathbf{x_{p}}}||_{2}
    \end{split}
\end{equation}
\normalsize

\paragraph{Multi-Objective Experiment}
Given the desirability of both Young's modulus and yield strength in a
 thermoset composite, a dual-objective Bayesian optimisation experiment was
 attempted, with the \textit{a posteriori} goal of sample-efficiently
 identifying the pareto front (\(\mathcal{P}\)) \cite{Garnett2023}.
The Pareto front consists of all samples which remain unsurpassed (by any
 other sample evaluated) in terms of at least one of their target variables
 (\(E\) or \(\sigma_{y}\)).
This Pareto front allows a rational agent to make subjective yet informed
 decisions about material composition in light of expected trade-offs in
 end properties.
Equation~\ref{eq:MOBOProblem} frames the multi-objective optimisation
 problem in terms of a dual-objective function \(\mathbf{f}\) (comprised of
 the single-objective functions \(f_{E}\) and \(f_{\sigma}\)) which returns
 a vector \(\mathbf{y}\) of the two targets (\(E\) and \(\sigma_{y}\)).
\begin{equation}\label{eq:MOBOProblem}
    \begin{split}
        \mathbf{y}&=\mathbf{f}(\mathbf{x})\\
        \{E,\sigma_{y}\}&=\{f_{E}(\mathbf{x}),f_{\sigma}(\mathbf{x})\}\\
    \end{split}
\end{equation}

An observed Pareto front may be obtained directly from a given dataset
 \(\mathcal{D}\) using the function \(f_\mathcal{P}\)
 (Eq.~\ref{eq:ObservedPFunc}).
\begin{equation}\label{eq:ObservedPFunc}
    \begin{split}
        \mathcal{P}&=f_\mathcal{P}(\mathcal{D})\\
        &=\{\mathbf{y}|\mathbf{y}\in\mathbf{Y},\nexists\:\mathbf{y}'\in\mathbf{Y}s.t.\:\mathbf{y}'\succ\mathbf{y}\}
    \end{split}
\end{equation}

Alternatively, where a multi-objective function \(\mathbf{f}\) is
 considered alongside a matrix of predictor variables \(\mathbf{X}\), the
 function \(f_\mathcal{P}^{*}\) may be used to produce a Pareto
 front (Eq.~\ref{eq:ParetoFrontFunc}).
\small
\begin{equation}\label{eq:ParetoFrontFunc}
    \begin{split}
        \mathcal{P}&=f_{\mathcal{P}}^{*}(\mathbf{X};\mathbf{f})\\
        &=\{\mathbf{f}(\mathbf{x})|\mathbf{x}\in\mathbf{X},\mathbf{f}(\mathbf{x})\succ\mathbf{f}(\mathbf{x}')\forall\mathbf{x}'\in\mathbf{X}\}
    \end{split}
\end{equation}
\normalsize

By retrieving all the raw mechanical test data associated with the samples
 evaluated in the three sequential experiments and their common prior, a
 56-sample dataset \(\mathcal{D}\) was generated comprised of a matrix of
 sampled predictors \(\mathbf{X}\) and a matrix of their respective
 targets \(\mathbf{Y}\) (\(\mathbf{E}\) and \(\bm{\sigma_{y}}\)).
Having conditioned two separate Gaussian processes on \(\mathbf{E}\) and
 \(\bm{\sigma_{y}}\) we obtained two separate posterior processes
 (\(\mathcal{GP}_{E}\) and \(\mathcal{GP}_{\sigma}\)); these
 provided initial prior processes for all subsequent sequential optimisation.
The pair of resulting processes were respectively described mathematically
 by their mean (\(m_{E}\), \(m_{\sigma}\)) and covariance functions
 (\(K_{E}\), \(K_{\sigma}\)) and computationally by their multivariate
 normal analogues (\(\bm{\mu}_{E}\), \(\bm{\mu}_{\sigma}\) and
 \(\bm{\Sigma}_{E}\), \(\bm{\Sigma}_{\sigma}\)).
Equation~\ref{eq:MOBOSampledFuncs} shows how a dual-objective surrogate
 function \(\tilde{\mathbf{f}}\) may be sampled from the pair of Gaussian
 processes (\(\mathcal{GP}_{E}\) and \(\mathcal{GP}_{\sigma}\)) previously
 conditioned.
\footnotesize
\begin{equation}\label{eq:MOBOSampledFuncs}
    \begin{split}
    \tilde{\mathbf{f}}\sim\:&\{\mathcal{GP}_{E},\mathcal{GP}_{\sigma}\}\\
    \{\tilde{f_{E}},\tilde{f_{\sigma}}\}\sim\:&\{\mathcal{GP}(m_{E},K_{E}),\mathcal{GP}(m_{\sigma},K_{\sigma})\}
    \end{split}
\end{equation}
\normalsize

Utility of a given Pareto front was evaluated using the hypervolume indicator
 utility function \(u_{hv}\) presented in Equation~\ref{eq:UtilityHV}, which
 is equal to the d-dimensional Lebesgue measure \(\lambda_{d}\) of the space
 dominated by a Pareto front \(\mathcal{P}\) and bounded from below by
 reference point \(r\) (\(r\) and \(v\) describe vertices).
\begin{equation}\label{eq:UtilityHV}
    u_{hv}(\mathcal{P};r)=\lambda_{d}(\bigcup_{v \in \mathcal{P}}[r,v])
\end{equation}

Equation~\ref{eq:FunctionHVI} shows how the utility function \(u_{hv}\) is
 used in the subsequent calculation of the one-step marginal gain known as
 the ``hypervolume improvement'' between the extant \(\mathcal{P}\) and the
 updated Pareto front \(\mathcal{P}'\) (which includes the putative point
 \(\mathbf{x}_{p}\)).
\begin{equation}\label{eq:FunctionHVI}
    f_{hvi}(\mathcal{P},\mathcal{P}';r)=u_{hv}(\mathcal{P}';r)-u_{hv}(\mathcal{P};r)
\end{equation}

The common multi-objective acquisition function by the name of ``expected
 hypervolume improvement'' is described by Equation~\ref{eq:AcqFuncEHVI},
 which takes a putative point \(\mathbf{x}_{p}\) and the current dataset
 \(\mathcal{D}\) and returns the instantaneous improvement metric.
Expectation manifests itself during the speculative update of the dataset
 \(\mathcal{D}\) to \(\mathcal{D}'\), where the two targets are drawn from
 their respective marginal Gaussian distributions (of their respectively
 updated posterior Gaussian processes) located at the putative point
 \(\mathbf{x}_p\); whose expectation (\(\mathbb{E}(\mathbf{x}_{p})\)) is
 naturally the mean of the marginal at that location
 (\(m_{E}(\mathbf{x}_{p})\) or \(m_{\sigma}(\mathbf{x}_{p})\)).
\small
\begin{equation}\label{eq:AcqFuncEHVI}
    \begin{split}
        &\alpha_{ehvi}(\mathbf{x}_{p};\mathcal{D})=f_{hvi}(f_{\mathcal{P}}(\mathcal{D}'),f_{\mathcal{P}}(\mathcal{D});r)\\
        &\text{where}\:\mathcal{D}'=\mathcal{D}\cup\{\mathbf{x}_{p},m_{E}(\mathbf{x}_{p}),m_{\sigma}(\mathbf{x}_{p})\}
    \end{split}
\end{equation}
\normalsize

This work adopted Daulton's re-formulation of the classic acquisition
 function (Eq.~\ref{eq:AcqFuncEHVI}) which takes into account the noisy
 nature of objective function evaluations \cite{Daulton2021}; although the
 true formulation contains an integral which is analytically
 intractable, it may be approximated using Monte Carlo integration.
Equation~\ref{eq:FunctionNEHVI} describes the integral-aprroximating 
 re-formulation of the acquisition function \(\hat{\alpha}\) which takes a
 putative point \(\mathbf{x}_{p}\) and the dataset \(\mathcal{D}\) and
 returns an instantaneous improvement value.
This ``noisy expected hypervolume improvement'' acquisition function
 iterates from \(1,\dots,N\) using the \(t\) counter, with each iteration
 drawing a random dual-objective function \(\tilde{\mathbf{f}}_{t}\) from
 the posterior gaussian process (see Eq.~\ref{eq:MOBOSampledFuncs}).
This function \(\tilde{\mathbf{f}}_{t}\) is used to speculatively produce
 Pareto fronts (\(\mathcal{P}\) and \(\mathcal{P}'\)) for both an original
 dataset \(\mathbf{X}\) and the updated dataset
 \(\mathbf{X}\cup\mathbf{x}_{p}\) respectively, before the hypervolume
 improvement is gauged using \(f_{hvi}\) and the many trials \(N\) averaged.
\small
\begin{equation}\label{eq:FunctionNEHVI}
    \begin{split}
        &\hat{\alpha}_{nehvi}(\mathbf{x}_{p};\mathbf{X})=\\
        &\frac{1}{N}\sum_{t=1}^{N}f_{hvi}(f_{\mathcal{P}}^{*}(\mathbf{X}\cup\mathbf{x}_{p};\tilde{\mathbf{f}}_{t}),f_{\mathcal{P}}^{*}(\mathbf{X};\tilde{\mathbf{f}}_{t});r)
    \end{split}
\end{equation}
\normalsize
This work practically applied the log implementation of the noise-aware
 \(\alpha_{nehvi}\) acquisition function
 ``qLogNoisyExpectedHypervolumeImprovement'' (``qLNEHVI'') which was shown to
 improve Bayesian optimisation performance in noisy settings
 \cite{Daulton2021}. 

\paragraph{Implementation}
All experiments were implemented in the Python programming language, with
 functionality spanning initial sampling techniques,
 adaptive sequential sampling methods, and an experimental workflow
 management system facilitating human-in-the-loop laboratory work; this
 package had the acronym ``BO4ACST'' (version 1.0.0)- ``Bayesian
 optimisation for atmospheric carbon sequestering thermosets''
 \cite{HojlundDodd2025}.
Note that Bayesian sequential optimisation techniques including Gaussian
 process regression, gaussian process conditioning, and sample utility
 evaluation were facilitated by the open-source ``adaptive experimentation
 platform'' or ``Ax'' (version 1.1.2) \cite{Facebook2025}.
Note that ``Ax'' makes use of a great deal of code from ``Botorch'',
 especially in it's implementation of custom Bayesian optimisation routines
 (e.g. non-standard acquisition function implementations)
 \cite{MetaPytorch2025B}.